% Bozza tolta dalla Sezione 3.1 (Panoramica): la distribuzione dei componenti avrà un capitolo a parte.
% Non è inclusa in main.tex; richiede le librerie TikZ già caricate (calc).

\subsection*{Distribuzione dei componenti}

I componenti non vengono eseguiti tutti nello stesso processo. Il sistema è distribuito su tre unità di
esecuzione, rappresentate nella Figura~\ref{fig:deployment}: due servizi containerizzati con Docker, il server di
inferenza vLLM e il database vettoriale Qdrant, e il processo Python dell'applicazione, che ospita l'interfaccia e
la logica della pipeline. Tutte e tre girano sulla stessa macchina e comunicano attraverso la rete locale, senza
alcuna chiamata verso servizi esterni.

\begin{figure}[H]
  \centering
  \begin{tikzpicture}[
      unit/.style={draw, thick, rounded corners=3pt, align=center, inner sep=2mm, font=\footnotesize},
      dock/.style={unit, fill=blue!6, text width=3.7cm},
      ext/.style={draw, rounded corners=2pt, align=center, inner sep=2mm, font=\footnotesize},
      link/.style={{Stealth[length=2mm]}-{Stealth[length=2mm]}, thick},
      lbl/.style={font=\scriptsize, align=center},
    ]
    \node[ext, text width=1.9cm] (browser) at (0,0) {Browser dello studente};
    \node[unit, text width=4.3cm] (app) at (4.9,0) {\textbf{Applicazione} (processo Python)\\[2pt]
      interfaccia Chainlit e pipeline\\ BGE-M3 e reranker (GPU~0)\\ docstore in memoria};
    \node[dock] (vllm) at (11.0,1.45) {\textbf{vLLM} (container Docker)\\[2pt]
      modello linguistico (GPU~1)\\ API compatibile OpenAI};
    \node[dock] (qdrant) at (11.0,-1.45) {\textbf{Qdrant} (container Docker)\\[2pt]
      vettori e metadati\\ volume persistente};
    \node[ext, text width=4.3cm] (offline) at (4.9,-3.1) {Script di indicizzazione e di valutazione
      (processi Python, su richiesta)};

    \draw[link] (browser) -- node[lbl, above] {HTTP,\\WebSocket} (app);
    \draw[link] (app.east) -- node[lbl, sloped, above] {HTTP} (vllm.west);
    \draw[link] (app.east) -- node[lbl, sloped, below] {gRPC} (qdrant.west);
    \draw[link, dashed] (offline.east) -| node[lbl, pos=0.75, left] {gRPC} (qdrant.south);
    \draw[link, dashed] (offline.east) -- (13.45,-3.1) |- node[lbl, pos=0.25, right] {HTTP} (vllm.east);
  \end{tikzpicture}
  \caption{Distribuzione del sistema su una singola macchina. vLLM e Qdrant sono servizi containerizzati;
  l'applicazione ospita l'interfaccia, la pipeline e i modelli di embedding e reranking. Gli script di
  indicizzazione e valutazione usano gli stessi servizi.}
  \label{fig:deployment}
\end{figure}

\paragraph{vLLM.} Il modello linguistico è servito da vLLM, un motore di inferenza progettato per servire modelli
linguistici con un elevato numero di richieste concorrenti. vLLM espone un'interfaccia HTTP compatibile con l'API
di OpenAI: l'applicazione invia richieste di tipo \emph{chat completion} e, per la risposta allo studente, ne
riceve i token in streaming man mano che vengono generati. Il modello viene pubblicato con un nome logico fisso,
indipendente dal modello effettivamente caricato: sostituire il modello, come è stato fatto per confrontarne
diversi nel Capitolo~\ref{cap:validazione}, richiede solo di cambiare un parametro di configurazione e riavviare il
container, senza modifiche al codice. Lo stesso servizio risponde a tutte le richieste rivolte al modello
linguistico, che per una singola domanda sono più d'una: la riscrittura della domanda (solo se la conversazione è
già avviata), la sua scomposizione e il riconoscimento del corso, eseguiti in parallelo, e infine la generazione
della risposta; nei casi in cui il recupero non trova nulla si aggiungono la ricerca con HyDE e il controllo delle
domande fuori tema. In fase di indicizzazione, lo stesso servizio descrive le immagini contenute nei documenti.
vLLM accorpa dinamicamente le richieste concorrenti in un'unica elaborazione sulla GPU (\emph{continuous
batching}), così che le chiamate parallele di una stessa domanda, o di più studenti contemporaneamente, non si
accodino una dopo l'altra.

\paragraph{Qdrant.} Qdrant conserva i vettori densi e sparsi dei frammenti insieme ai loro metadati, su cui sono
definiti indici per filtrare la ricerca, ad esempio per corso di laurea. L'applicazione lo interroga via gRPC: per
ogni (sotto-)domanda invia un'unica richiesta che esegue in parallelo la ricerca densa e quella sparsa e ne fonde i
risultati direttamente nel database, restituendo i candidati già ordinati. I dati sono salvati in un volume Docker
persistente, che sopravvive al riavvio del container.

\paragraph{Applicazione.} Il processo dell'applicazione ospita l'interfaccia Chainlit, raggiungibile dal browser
via HTTP, con una connessione WebSocket su cui viaggiano i messaggi della chat e la risposta in streaming. Nello
stesso processo vengono eseguiti i modelli di embedding e di reranking, sulla GPU non occupata da vLLM: sono modelli
molto più piccoli del modello linguistico e vengono invocati come funzioni locali, senza passare per la rete. Il
docstore, che contiene il testo dei frammenti e le sezioni usate dalla fusione gerarchica, è un file su disco
caricato in memoria all'avvio. Gli script di indicizzazione e di valutazione sono processi Python separati, eseguiti
su richiesta, che usano gli stessi servizi vLLM e Qdrant.

La configurazione dell'intero sistema --- indirizzi dei servizi, GPU assegnate, modello caricato in vLLM, nome della
collezione in Qdrant --- è raccolta in un unico file di variabili d'ambiente, letto sia da Docker sia
dall'applicazione. La Sezione~\ref{sec:infrastruttura} approfondisce lo stack tecnologico e le considerazioni sul
deploy.
